\part{Pourquoi et comment compresser une image}
\section{Réversibilité et théorie de l'information}
\subsection{Trois pertes à ne pas confondre}
La conversion analogique-numérique quantifie une mesure physique.
Un codec avec pertes peut ensuite supprimer de l'information
numérique. Enfin, un rendu peut écrêter le gamut ou les hautes
lumières. Une compression \emph{sans perte} garantit seulement
la reconstruction exacte des données auxquelles elle s'applique,
pas la préservation de toute la scène d'origine.

Si un encodeur $E$ et un décodeur $D$ sont sans perte sur un domaine
$\mathcal X$, on exige $D(E(x))=x$ pour tout $x\in\mathcal X$.
Une transformation inversible est une bijection sur ce domaine.
L'arrondi de 12,4 et de 12,49 vers 12 est au contraire plusieurs-à-un :
l'inverse exact est impossible sans information supplémentaire.

\subsection{Entropie : le coût moyen d'une surprise}
Si un symbole $i$ apparaît avec probabilité $p_i$, l'entropie est
\begin{equation}
H=-\sum_i p_i\log_2p_i.
\end{equation}
Un symbole certain apporte zéro bit de surprise ; deux symboles
équiprobables apportent un bit. Cette lecture probabiliste de
l'information vient de Shannon~\cite{shannon}. Elle ne mesure
ni la beauté de l'image ni sa valeur scientifique.

Exemple original : les probabilités $(1/2,1/4,1/8,1/8)$ donnent
$H=1{,}75$ bit par symbole. Un stockage fixe à deux bits consomme
deux bits par symbole, tandis qu'un code préfixe adapté peut
atteindre ici exactement 1,75 bit en moyenne. Le fichier réel
ajoute une description du code, des en-têtes et des limites de blocs.

Il est impossible de raccourcir sans perte \emph{tous} les fichiers :
il y a $2^n$ suites de $n$ bits et seulement $2^n-1$ suites de
longueur strictement inférieure à $n$. Une compression efficace
exploite donc des régularités de la famille d'images rencontrée.
Un bruit aléatoire indépendant se compresse souvent moins bien
qu'une plage uniforme, même si les deux ont les mêmes dimensions.

\section{Prédiction, différences et codes entropiques}
\subsection{Prédire pour rendre les données plus faciles à coder}
Si une ligne vaut $(100,102,101,103)$, conserver le premier nombre
et les différences donne $(100,2,-1,2)$. La transformation est
réversible : les sommes cumulées restituent la ligne.
Le nombre de valeurs n'a pas changé, mais les résidus se
concentrent autour de zéro et peuvent coûter moins de bits.

Avec un prédicteur $P_i$ calculé à partir de valeurs déjà connues :
\begin{equation}
r_i=x_i-P_i,\qquad x_i=P_i+r_i.
\end{equation}
Le décodeur doit employer exactement le même prédicteur et le
même ordre de parcours. Un prédicteur adaptatif doit préciser
son état initial et ses règles de mise à jour.

Une soustraction peut aussi être effectuée modulo 256 sur des
octets. Par exemple $3-250\equiv9\pmod{256}$, puis
$250+9\equiv3\pmod{256}$. Le résidu stocké est 9, pas un entier
signé de valeur $-247$. Cette convention permet des filtres
réversibles sans changer la largeur des octets.

\subsection{RLE : compter les répétitions}
La suite \texttt{AAAAABBBCC} peut être représentée par
$(5,A),(3,B),(2,C)$. Sur une photographie bruitée, les suites
constantes sont courtes ; ce codage peut même grossir les données.
Des marqueurs doivent distinguer une répétition d'une séquence
littérale. Le principe est simple, ses conventions binaires
exactes appartiennent au format qui l'utilise.

\subsection{Huffman : des mots courts pour les symboles fréquents}
Dans notre exemple de probabilités, on peut choisir
$A\mapsto0$, $B\mapsto10$, $C\mapsto110$, $D\mapsto111$.
Aucun mot n'est le début d'un autre : le code est préfixe.
La suite $ABCD$ devient \texttt{010110111} et se relit sans
séparateurs en parcourant un arbre. La longueur moyenne vaut
$0{,}5\times1+0{,}25\times2+0{,}125\times3+0{,}125\times3=1{,}75$.

Huffman code des symboles discrets avec des longueurs entières.
Un codage arithmétique ou par intervalles représente une suite
entière par une sous-plage, permettant une adaptation plus fine
aux probabilités. Les implémentations doivent définir les
arrondis et la renormalisation pour rester synchronisées.

\subsection{LZ77 et DEFLATE}
Un dictionnaire glissant remplace une séquence répétée par une
distance en arrière et une longueur. Il ne cherche pas des
objets visuels : il cherche des octets déjà rencontrés.
DEFLATE combine des références de ce type et un codage Huffman ;
son format est défini par la RFC~1951~\cite{deflate}.
Zlib ajoute un encadrement et un contrôle d'intégrité décrits
par la RFC~1950~\cite{zlib}. Zlib, DEFLATE et ZIP sont donc
trois notions apparentées, mais non interchangeables.

\subsection{Rice/Golomb : encoder des petits résidus}
Pour un entier non négatif $u$ et un paramètre $k$, un code de
Rice sépare quotient $q=\lfloor u/2^k\rfloor$ et reste
$r=u\bmod2^k$. Le quotient est écrit en unaire, le reste sur
$k$ bits. Avec la convention « $q$ uns suivis d'un zéro »,
$u=13$, $k=2$ donne $q=3$, $r=1$, donc \texttt{1110 01}.

Pour transformer un résidu signé en entier non négatif, on peut
utiliser $u=2r$ si $r\ge0$, sinon $u=-2r-1$. Les valeurs
$0,-1,1,-2,2$ deviennent $0,1,2,3,4$. Le choix de $k$ et les
conventions varient selon le codec ; cet exemple n'est pas un
décodeur Canon, Sony ou Nikon.

\section{Transformées : changer de coordonnées}
\subsection{Pourquoi passer aux fréquences ?}
Un bloc uniforme contient beaucoup de valeurs égales. Dans une
base adaptée, il peut être représenté par un coefficient moyen
et des détails nuls. Une transformée inversible ne compresse
pas nécessairement à elle seule ; elle concentre l'information
pour une quantification et un codage plus efficaces.

Pour une transformation orthonormale $A$, $c=Ax$ et
$x=A^{\mathsf T}c$. L'énergie quadratique se conserve :
$\|c\|_2^2=\|x\|_2^2$. Les erreurs dans les coefficients se
traduisent alors directement en erreurs quadratiques dans
l'image, sous réserve d'arrondis de calcul.

\subsection{DCT 2D et exemple JPEG}
Une DCT orthonormale d'un bloc $8\times8$ centré est
\begin{equation}
F(u,v)=\frac14\alpha(u)\alpha(v)
\sum_{x=0}^{7}\sum_{y=0}^{7}f(x,y)
\cos\!\frac{(2x+1)u\pi}{16}
\cos\!\frac{(2y+1)v\pi}{16},
\end{equation}
avec $\alpha(0)=1/\sqrt2$, $\alpha(k)=1$ sinon.
Le coefficient $(0,0)$ est le DC ; les autres décrivent des
variations spatiales AC. Dans un exemple 8 bits classique,
on soustrait 128 avant la transformée.

Si tous les pixels valent 140, le bloc centré vaut 12 ;
$F(0,0)=96$ et les autres coefficients sont nuls en calcul exact.
L'inverse retrouve 12 partout, puis 140 après réaddition de 128.
Cette construction mathématique explique la concentration de
l'énergie ; le cadre normatif JPEG est T.81~\cite{jpeg}.

\subsection{Quantification : l'étape irréversible}
On calcule pour chaque coefficient
\begin{equation}
\widehat F(u,v)=\operatorname{round}\!\frac{F(u,v)}{Q(u,v)},
\qquad F_{\mathrm{reconstruit}}=Q\widehat F.
\end{equation}
Avec $F=23$ et $Q=10$, on stocke 2 et on reconstruit 20.
Les valeurs de 15 à moins de 25, selon la convention d'arrondi,
peuvent partager ce code ; l'information est perdue.
Les hautes fréquences sont souvent quantifiées plus fortement.
Des détails faibles deviennent nuls, ce qui favorise ensuite
les suites de zéros et les codes entropiques.

La « qualité JPEG 90 » est un paramètre d'encodeur, pas une
unité normalisée garantissant la même table entre bibliothèques.
Réexporter plusieurs fois avec pertes peut accumuler des erreurs,
même si certaines configurations se stabilisent partiellement.

\subsection{Zigzag, DC différentiel et sous-échantillonnage}
Un parcours zigzag regroupe approximativement les coefficients
des basses vers les hautes fréquences. Les zéros consécutifs
peuvent être codés efficacement. Le DC est souvent prédit à
partir du bloc précédent ; les marqueurs de redémarrage
réinitialisent des états définis par le codec.

Avant le codage, une image peut être convertie en une composante
de luma et deux de chroma. Le format 4:2:0 usuel conserve une
grille de luma et deux grilles de chroma à demi-largeur et
demi-hauteur. Pour un rectangle de dimensions paires, cela
représente environ $1{,}5WH$ échantillons contre $3WH$ en 4:4:4.
L'emplacement exact des échantillons de chroma et les règles de
reconstruction comptent. Les textes colorés et contours
saturés peuvent souffrir davantage que les surfaces photographiques.

\subsection{Ondelettes et lifting}
Les ondelettes décomposent l'image par échelles plutôt qu'en
blocs DCT indépendants. Une étape de Haar produit, pour deux
valeurs, $(a,d)=((x_0+x_1)/\sqrt2,(x_0-x_1)/\sqrt2)$.
Appliquée sur lignes puis colonnes, elle donne un fond et
trois orientations de détails. Le cours 10 développe cette
construction et son lien avec l'estimation du bruit.

Le lifting construit des transformées par prédiction et mise
à jour. Avec des règles d'arrondi entières symétriques,
certaines variantes sont exactement réversibles. Quantifier
leurs sous-bandes produit en revanche un codec avec pertes.
Dire « ondelettes » ne suffit donc pas à conclure « sans perte ».

\section{Ce qu'un décodeur doit reconstruire}
Un décodeur JPEG à DCT ne lit pas directement les pixels : il
parse les marqueurs, reconstruit les tables, décode les
symboles, remet les coefficients dans les blocs, déquantifie,
effectue la DCT inverse, reconstruit les composantes et applique
les conversions demandées. Une décompression prédictive suit
une chaîne différente. Les détails normatifs fixent la
précision, les modes et les conventions.

\resultat{Décompresser exactement un fichier avec pertes ne restaure
pas les valeurs originales : cela restaure exactement, ou suivant
la précision autorisée du décodeur, la représentation déjà appauvrie
que l'encodeur a choisie.}
